\paragraph{Model merging: constructions without limits.}
A wave of post-hoc model-merging methods has emerged since Task
Arithmetic~\cite{ilharco2023task} showed that adding fine-tuning
deltas of separately-trained task vectors recovers non-trivial
multi-task performance. Fisher-weighted averaging~\cite{matena2022fisher},
sparsification-based approaches (TIES~\cite{yadav2023ties},
DARE~\cite{yu2024dare}), and more recent rotation- and
subspace-alignment families (KnOTS~\cite{stoica2025knots},
TSPA~\cite{tspa2025}, DO-Merging~\cite{domerging2025},
Core~Space~\cite{panariello2025core}, TARA~\cite{tara2025},
ARM~\cite{arm2026streaming}) all seek to reduce destructive
interference between task vectors before averaging. A parallel
line quantizes task vectors at merge time to shrink the storage
footprint: Task-Vector Quantization (TVQ)~\cite{kim2025tvq} and
1-bit-Merging~\cite{bit1merging2025} achieve competitive accuracy
at one to four bits per coordinate. What unifies these methods is
that they are \emph{constructions}: each proposes a particular
encoder and evaluates it empirically. None establishes a
\emph{fundamental limit} on merging quality at a given storage or
rate budget. This paper supplies that limit, and — in the generic
regime — a matching construction.

\paragraph{Existing theoretical analyses of merging.}
The closest theoretical work on merging falls into three strands.
\emph{Weight disentanglement} (Ortiz-Jimenez et
al.~\cite{ortizjimenez2023disentanglement}) identifies the
rank-$r$ Hessian structure of LoRA fine-tunes and relates task
arithmetic's empirical success to locally-additive task
functionals; this provides the $H_t = P_{V_t}$ (and more
generally $H_t = U_t D_t U_t^\top$) formalism we adopt in
\S\ref{sec:setup}, but does not give rate-distortion bounds.
\emph{Task-Vector Bases}~\cite{jang2025taskvectorbases} posits
that fine-tuned weights lie in a thin Gaussian shell — a
geometric assumption, not information-theoretic. \emph{On Task
Vectors and Gradients}~\cite{gradients2025taskvectorsgradients}
interprets task vectors as single-step gradient updates and
provides a second-order Taylor-expansion bound on task
arithmetic error. ATM~\cite{atm2024alternating} takes the
noisy-GD interpretation further. None of these works frame
merging as lossy source coding or derive Shannon-style limits.
Our work is the first to pose merging as a multi-source
rate-distortion problem with worst-task distortion as the figure
of merit.

\paragraph{Rate-distortion and rotation-for-quantization ancestry.}
The information-theoretic toolkit we use is standard. Classical
rate-distortion theory~\cite{shannon1959rd,coverthomas} gives
the single-source setting; multi-description coding
\cite{elgamalcover1982md} addresses multi-receiver lossy coding
but under average — not maximum — distortion, which precludes
direct application to merging (Lemma~\ref{lem:id} shows the
max-distortion formulation is strictly stronger and exhibits a
different RD function). The direct intellectual ancestor of our
achievability construction is
TurboQuant~\cite{zandieh2025turboquant}, which establishes
Shannon-style lower bounds on randomized vector quantization via
Yao's minimax principle and achieves them up to a universal
constant via Haar-rotation + uniform scalar quantization. Our
Theorem~\ref{thm:achv} reduces to TurboQuant Theorem 3 in the
$T=1$ single-task limit. A further line of work uses structured
rotations (Hadamard / Walsh-Hadamard) as pre-conditioners for
quantization in the single-model setting:
QuIP~\cite{tseng2024quip}, QuaRot~\cite{ashkboos2024quarot},
SpinQuant~\cite{liu2025spinquant}, and
RoLoRA~\cite{xi2024rolora}. These motivate our choice of
randomized orthogonal mixers in the achievability algorithm but
do not address the multi-source / merging setting.

\paragraph{The November 2025 negative result, and why a lower
bound still matters.} A recent large-scale empirical study
\cite{systematic2025merging} reports that across four LLM
families and sixteen benchmarks, only Task Arithmetic reliably
produces constructive interference; subspace-alignment methods
(KnOTS-style) often underperform or fail. This might seem to
undermine the case for theoretical work on merging. It does not.
Their result is an \emph{upper bound} on specific algorithms —
evidence that those constructions are suboptimal at LLM scale.
Our Theorem~\ref{thm:general} is a \emph{lower bound} on every
algorithm — it says that whatever the best merging procedure is,
it cannot beat $B^2 f_{\mathrm{floor}} + c_{\mathrm{TQ}}
(B^2 \deff / (Tr)) \cdot 2^{-2R/\deff}$. The negative result and
our lower bound are compatible and mutually reinforcing, and our
own measurements (\S\ref{sec:exp}) pin down the mechanism. We
find that real LoRA task subspaces are linearly
\emph{independent} ($\deff = Tr$ in every layer of all four base
models we test), i.e.\ the floor-zero regime. In that regime a
subspace-alignment method has no overlap to exploit, so it cannot
improve on naive averaging — and indeed we measure KnOTS to be
statistically indistinguishable from Task Arithmetic, while only
TIES (which resolves sign-level interference rather than subspace
overlap) helps. This explains the
\cite{systematic2025merging} negative result not by an
irreducible floor but by its opposite: with no subspace overlap,
the floor is zero and the entire merging gap is algorithmic
slack above a beatable target. The lower bound's
\emph{floor-positive} regime — substantial subspace overlap, no
amount of rate sufficient — is theoretically characterized here
but did not arise empirically for the task families we tried;
identifying real task sets that induce it, and fully
characterizing achievability there, remain open
(Remark~\ref{rem:sharedv-gap}).
